\documentclass[11pt,a4paper]{article}

% ==================================================
% IDIOMA Y TIPOGRAFÍA
% ==================================================

\usepackage[utf8]{inputenc}
\usepackage[T1]{fontenc}
\usepackage[spanish]{babel}
\usepackage{lmodern}
\usepackage{microtype}

% ==================================================
% DISEÑO DE PÁGINA
% ==================================================

\usepackage[left=1.7cm,right=1.7cm,top=1.5cm,bottom=1.4cm]{geometry}
\usepackage{enumitem}
\usepackage[hidelinks]{hyperref}

\setlength{\parindent}{0pt}
\setlength{\parskip}{2pt}
\pagestyle{empty}

% ==================================================
% COMANDOS PARA MANTENER UN FORMATO UNIFORME
% ==================================================

\newcommand{\seccion}[1]{%
  \vspace{8pt}
  \begin{center}\textbf{#1}\end{center}
  \vspace{-5pt}\hrule\vspace{7pt}
}

\newcommand{\cargo}[4]{%
  \textbf{#1}\hfill #2\\
  \textbf{#3}\hfill \textit{#4}\par\vspace{2pt}
}

\newenvironment{funciones}{%
  \begin{itemize}[
    leftmargin=1.1em,
    itemsep=1pt,
    topsep=2pt,
    parsep=0pt
  ]
}{\end{itemize}\vspace{3pt}}

% ==================================================
% INICIO DEL DOCUMENTO
% ==================================================

\begin{document}

% ==================================================
% DATOS PERSONALES Y CONTACTO
% ==================================================

\begin{center}

  {\Large\textbf{Daniel Leandro Pajita Dominguez}}\\[4pt]

  Lima, Perú \textbullet\
  \href{mailto:danielleandro.pd@gmail.com}
  {danielleandro.pd@gmail.com}\\[3pt]

  \href{https://danielpajita.github.io/}{Página personal}
  \textbullet\
  \href{https://github.com/DanielPajita}{GitHub}
  \textbullet\
  \href{https://www.linkedin.com/in/daniel-leandro-pajita-dominguez-93657a265/}{LinkedIn}\\[4pt]

  \hrulefill

\end{center}

% ==================================================
% EDUCACIÓN
% ==================================================

\seccion{Educación}

\textbf{Pontificia Universidad Católica de Chile (UC)}
\hfill Chile\\
Maestría en Ciencia de Datos (en curso)
\hfill 2025--2027\\[5pt]

\textbf{Universidad Nacional del Callao (UNAC)}
\hfill Perú\\
Bachiller en Economía
\hfill 2015

% ==================================================
% EXPERIENCIA PROFESIONAL
% ==================================================

\seccion{Experiencia profesional}

% PRODUCE

\cargo
  {Ministerio de la Producción (PRODUCE)}
  {Lima, Perú}
  {Consultor económico}
  {Oct. 2024 -- actualmente}

\begin{funciones}

  \item Análisis económico y evaluación de instrumentos de apoyo a empresas; procesamiento de registros administrativos y construcción de indicadores.

  \item Elaboración de informes técnicos para el seguimiento y la evaluación de intervenciones públicas.

\end{funciones}

% SUNEDU

\cargo
  {Superintendencia Nacional de Educación Superior Universitaria (SUNEDU)}
  {Lima, Perú}
  {Analista de datos}
  {Oct. 2020 -- Jul. 2024}

\begin{funciones}

  \item Elaboración de estudios, reportes estadísticos y boletines sobre el sistema universitario peruano.

  \item Procesamiento de registros administrativos reportados mediante TUNI y análisis de matrícula e inserción laboral.

\end{funciones}

% GRADE

\cargo
  {Grupo de Análisis para el Desarrollo (GRADE)}
  {Lima, Perú}
  {Asistente de investigación}
  {Mar. 2020 -- Ago. 2020}

\begin{funciones}

  \item Procesamiento de encuestas de hogares y empleo para estudios de desarrollo humano y políticas laborales durante la pandemia de COVID-19.

\end{funciones}

% MTPE

\cargo
  {Ministerio de Trabajo y Promoción del Empleo (MTPE)}
  {Lima, Perú}
  {Analista de mercado de trabajo}
  {Oct. 2017 -- Feb. 2020}

\begin{funciones}

  \item Elaboración de informes nacionales sobre empleo, juventud y demanda laboral; análisis econométrico y asistencia técnica.

  \item Aplicación de modelos de duración para estudiar cambios en la remuneración mínima y el empleo formal.

\end{funciones}

% ARELLANO MARKETING

\cargo
  {Arellano Marketing}
  {Lima, Perú}
  {Analista de macrodata}
  {Jul. 2016 -- Sep. 2017}

\begin{funciones}

  \item Análisis de información económica y sociodemográfica de fuentes secundarias; tratamiento de bases de datos y revisión de metodologías cuantitativas.

\end{funciones}

% ==================================================
% COMPETENCIAS TÉCNICAS
% ==================================================

\seccion{Competencias técnicas}

\textbf{Programación y software:}
Stata, Python, R, SQL y Git/GitHub.\\
\textbf{Métodos:}
econometría aplicada, evaluación de impacto, análisis de encuestas, modelos de duración y aprendizaje automático.

% Añadir el apartado de idiomas cuando se defina
% el nivel de dominio correspondiente.

% ==================================================
% PUBLICACIONES Y TRABAJOS
% ==================================================

\seccion{Publicaciones y trabajos}

\begin{itemize}[
  leftmargin=1.1em,
  itemsep=3pt,
  topsep=0pt
]

  \item \textbf{Pajita Dominguez, D.} y Baldeón Arrunátegui, A. (2017).
  ``Sistema de demanda QUAIDS para bebidas no alcohólicas en hogares pobres y no pobres del Perú, 2012--2016''.
  \textit{Horizonte Económico}, núm. 6.

  \item Huarancca Bellido, M. y \textbf{Pajita Dominguez, D.} (2017).
  ``Determinantes del embarazo y maternidad en los adolescentes del Perú''.
  \textit{Horizonte Económico}, núm. 5, pp. 113--126.

  \item Malo Serrano, M., Castillo, N. y \textbf{Pajita, D.} (2017).
  ``La obesidad en el mundo''.
  \textit{Anales de la Facultad de Medicina}, 78(2), pp. 67--72.

\end{itemize}

\vspace{3pt}

\href{https://danielpajita.github.io/publicaciones/}
{\underline{Consultar listado completo de publicaciones, informes y colaboraciones.}}

% ==================================================
% FORMACIÓN COMPLEMENTARIA
% ==================================================

\seccion{Formación complementaria}

\begin{itemize}[
  leftmargin=1.1em,
  itemsep=1pt,
  topsep=0pt
]

  \item Evaluación de Impacto de Programas de Desarrollo, GRADE (60 horas).

  \item Diploma en Econometría, Grupo Iddea (96 horas).

  \item Econometría Intermedia, Grupo Lambda (64 horas).

  \item Manejo de Bases de Datos y Econometría Aplicada con RStudio, Económica (25 horas).

  \item Stata para Economistas, Grupo Lambda (24 horas).

\end{itemize}

% ==================================================
% FIN DEL DOCUMENTO
% ==================================================

\end{document}